% rtex-serve-patches.tex: kernel/package shortcuts installed in the fast server only, after
% \begin{document} and before the serve loop starts. Each one removes work whose result can
% never be observed in a unit box (running-head marks, log-only messages) or memoizes a
% deterministic computation. Every patch checks that the macro it replaces still has the
% definition it was written against and is skipped otherwise. Row-exactness against the
% background pass is checked by `rtex verify` (docs/correctness.md).
\makeatletter
%
% 1. Running-head marks. \markright/\markboth insert \marks nodes that only the output routine
%    reads; the unit box is never shipped. Keep the one typesetting side effect (a \nobreak
%    penalty in vertical mode after a heading).
\ifcsname markright \endcsname
  \expandafter\def\csname markright \endcsname#1{\if@nobreak\ifvmode\nobreak\fi\fi}%
  \expandafter\def\csname markboth \endcsname#1#2{\if@nobreak\ifvmode\nobreak\fi\fi}%
\fi
%
% 2. Math font setup. \glb@settings rebuilds \math@fonts (3 \pickup@font per math group, ~40
%    per call, each through microtype's font hook) every time math is entered at a size other
%    than the last one, e.g. for every footnote mark. The result only depends on the math
%    version, the size and the \mv@<version> list, so it is memoized per (version, size) and
%    reused while \mv@<version> is unchanged (\ifx on the token list).
\def\rtex@glb@settings@expected{%
     \expandafter\ifx\csname S@\f@size\endcsname\relax
       \calculate@math@sizes
     \fi
     \csname S@\f@size\endcsname
     \ifmath@fonts
       \begingroup
         \escapechar\m@ne
         \csname mv@\math@version \endcsname
         \globaldefs\@ne
         \math@fonts
         \let \glb@currsize \f@size
       \endgroup
        \the\every@math@size
     \fi
}
\ifx\rtex@glb@settings@expected\glb@settings
  \def\glb@settings{%
     \expandafter\ifx\csname S@\f@size\endcsname\relax
       \calculate@math@sizes
     \fi
     \csname S@\f@size\endcsname
     \ifmath@fonts
       \begingroup
         \escapechar\m@ne
         \expandafter\ifx\csname rtex@mv@\math@version/\f@size\expandafter\endcsname
                         \csname mv@\math@version\endcsname
           \expandafter\let\expandafter\math@fonts\csname rtex@mf@\math@version/\f@size\endcsname
         \else
           \csname mv@\math@version \endcsname
           \global\expandafter\let\csname rtex@mf@\math@version/\f@size\endcsname\math@fonts
           \global\expandafter\let\csname rtex@mv@\math@version/\f@size\expandafter\endcsname
                                  \csname mv@\math@version\endcsname
         \fi
         \globaldefs\@ne
         \math@fonts
         \let \glb@currsize \f@size
       \endgroup
        \the\every@math@size
     \fi
  }%
\fi
%
% 3. graphics: the two file-existence probes per \includegraphics (\IfFileExists in
%    \Gin@getbase, \openin in \Gread@pdftex) are remembered for files that were found; the
%    image resource itself is already cached by luatex.def. Log-only messages are dropped.
\IfPackageLoadedTF{graphics}{%
  \let\Gin@log\@gobble
  \ifcsname Gread@@pdftex\endcsname
    \let\rtex@Gread@pdftex\Gread@pdftex
    \def\Gread@pdftex#1{%
      \ifcsname rtex@seen@#1\endcsname
        \Gread@@pdftex{#1}%
      \else
        \rtex@Gread@pdftex{#1}%
        \ifcsname #1 image\Gin@attr@hash\endcsname
          \global\expandafter\let\csname rtex@seen@#1\endcsname\@empty
        \fi
      \fi}%
  \fi
  \def\rtex@Gin@getbase@expected#1{%
    \edef\Gin@tempa{%
      \def\noexpand\@tempa####1#1\space{%
        \def\noexpand\Gin@base{####1}}}%
    \IfFileExists{\filename@area\filename@base#1}%
      {\Gin@tempa
       \edef\uq@filef@und{\expandafter\unquote@name
                          \expandafter{\@filef@und}}%
       \expandafter\@tempa\uq@filef@und
       \edef\Gin@ext{#1}}{}}%
  \ifx\rtex@Gin@getbase@expected\Gin@getbase
    \def\rtex@IfFileExists#1#2#3{%
      \ifcsname rtex@ff@#1\endcsname
        \edef\@filef@und{\csname rtex@ff@#1\endcsname}%
        \expandafter\rtex@IfFileExists@hit
      \else
        \expandafter\rtex@IfFileExists@i
      \fi{#2}{#3}{#1}}%
    \def\rtex@IfFileExists@hit#1#2#3{#1}%
    \def\rtex@IfFileExists@i#1#2#3{%
      \IfFileExists{#3}%
        {\global\expandafter\let\csname rtex@ff@#3\expandafter\endcsname\csname @filef@und\endcsname#1}%
        {#2}}%
    \def\Gin@getbase#1{%
      \edef\Gin@tempa{%
        \def\noexpand\@tempa####1#1\space{%
          \def\noexpand\Gin@base{####1}}}%
      \rtex@IfFileExists{\filename@area\filename@base#1}%
        {\Gin@tempa
         \edef\uq@filef@und{\expandafter\unquote@name
                            \expandafter{\@filef@und}}%
         \expandafter\@tempa\uq@filef@und
         \edef\Gin@ext{#1}}{}}%
  \fi
}{}
%
% 4. Package/class info messages only reach the log, which the server does not keep.
\let\GenericInfo\@gobbletwo
%
% 5. Title block. article/report (no titlepage) build it as a center environment, then make
%    \maketitle, \title, \author, \date, \thanks and \and unusable (\global\let ...\relax) and
%    empty \@title & co. The server must typeset the block again after every edit, so the
%    kernel definitions are saved here and restored after each compile (\rtex@restoretitle,
%    appended to the compile's tail tokens). \@maketitle loses its page-level material
%    (\newpage, \null, the \vskips outside the center environment): the capture's unit is the
%    environment, and the vertical glue around it comes from the placements.
\def\rtex@maketitle@expected{%
  \newpage
  \null
  \vskip 2em%
  \begin{center}%
  \let \footnote \thanks
    {\LARGE \@title \par}%
    \vskip 1.5em%
    {\large
      \lineskip .5em%
      \begin{tabular}[t]{c}%
        \@author
      \end{tabular}\par}%
    \vskip 1em%
    {\large \@date}%
  \end{center}%
  \par
  \vskip 1.5em}
\ifx\rtex@maketitle@expected\@maketitle
  \def\@maketitle{%
    \begin{center}%
    \let \footnote \thanks
      {\LARGE \@title \par}%
      \vskip 1.5em%
      {\large
        \lineskip .5em%
        \begin{tabular}[t]{c}%
          \@author
        \end{tabular}\par}%
      \vskip 1em%
      {\large \@date}%
    \end{center}%
    \par}%
  \let\rtex@save@maketitle\maketitle
  \let\rtex@save@@maketitle\@maketitle
  \let\rtex@save@title\title
  \let\rtex@save@author\author
  \let\rtex@save@date\date
  \let\rtex@save@thanks\thanks
  \let\rtex@save@and\and
  \let\rtex@save@@title\@title
  \let\rtex@save@@author\@author
  \let\rtex@save@@date\@date
  \let\rtex@save@@thanks\@thanks
  \edef\rtex@save@topnum{\the\@topnum}%
  \def\rtex@restoretitle{%
    \global\let\maketitle\rtex@save@maketitle
    \global\let\@maketitle\rtex@save@@maketitle
    \global\let\title\rtex@save@title
    \global\let\author\rtex@save@author
    \global\let\date\rtex@save@date
    \global\let\thanks\rtex@save@thanks
    \global\let\and\rtex@save@and
    \global\let\@title\rtex@save@@title
    \global\let\@author\rtex@save@@author
    \global\let\@date\rtex@save@@date
    \global\let\@thanks\rtex@save@@thanks
    \global\@topnum\rtex@save@topnum\relax
    \global\@specialpagefalse}%
\fi
%
\makeatother
